%% Applied Computing and Geosciences (Elsevier)
%% Build PDF: pdflatex -> bibtex -> pdflatex -> pdflatex
%% Build DOCX: ./make_docx.sh
\documentclass[preprint,12pt,authoryear]{elsarticle}

\usepackage{amsmath,amssymb}
\usepackage{graphicx}
\usepackage{booktabs}
\usepackage{multirow}
\usepackage[T1]{fontenc}
\usepackage[utf8]{inputenc}
\usepackage[hidelinks,bookmarks=false,pdfpagemode=UseNone]{hyperref}
\hypersetup{colorlinks=false,pdfborder={0 0 0}}

\journal{Applied Computing and Geosciences}

\newcommand{\Mgha}{Mg\,ha\textsuperscript{-1}}
\newcommand{\eqnum}[1]{\qquad\qquad\mathrm{(#1)}}

\begin{document}

\begin{frontmatter}

\title{Geostatistics belongs in the hyperparameter space, and coverage
belongs in the design-based assessment}

\author[uf]{Hugo Rodrigues}
\ead{rodrigues.machado.hugo@gmail.com}

\address[uf]{University of Florida, Range Cattle Research and Education
Center, 3401 Experiment Station, Ona, FL 33865, USA}

\begin{abstract}
Map accuracy must be estimated from a probability sample by design-based
inference, not by cross-validation. That argument settled how the point map
should be judged, but it left two questions open. Here we address both. First,
we argue that the geostatistical decisions still made by habit in digital soil
mapping, the variogram family, its range, the lag binning, the estimator and
the anisotropy, are hyperparameters and should be selected out of fold like any
other. We are explicit that geostatistics then supplies the hypothesis and
feature space rather than the final estimator: on the Amazonian above-ground
biomass data of \citet{wadoux2021} the corrector the data selects is a residual
forest in 68 of 90 repetitions and ordinary kriging in three, yet the resulting
estimator lowers design-based RMSE under all three sampling designs (32.49 to
31.62, 32.28 to 31.40 and 35.78 to 34.54 \Mgha{}), winning 28, 30 and 29 of 30
paired repetitions. The conventional variogram recipe is displaced in most
repetitions, and the procedure discriminates: it refuses an anisotropic geometry
under clustered sampling, where the design cannot support estimating direction,
while replacing the lag binning most often in exactly that design. Second, we
argue that prediction-interval coverage is a population parameter no different
from RMSE, so the same design-unbiased estimator applies. This matters because
it exposes a failure that is otherwise invisible: cross-fitted conformal
intervals are calibrated and 17 per cent sharper than quantile regression
forests under representative sampling, but under clustered sampling they
under-cover at 0.84 against a nominal 0.90, and only a spatial-block
calibration restores validity. Coverage, like accuracy, must be assessed
design-based.
\end{abstract}

\begin{keyword}
digital soil mapping \sep design-based inference \sep prediction intervals \sep
conformal prediction \sep variogram \sep random forest
\end{keyword}

\end{frontmatter}

\section{Introduction}

\citet{wadoux2021} showed that map accuracy indices must be estimated from a
probability sample by design-based inference, and that spatial cross-validation
has no theoretical underpinning for this purpose. The argument is now widely
accepted. It concerns, however, only the point map and only its accuracy, and
it leaves two questions that we take up here.

The first is where geostatistics belongs. Digital soil mapping has largely
replaced kriging with machine learning \citep{mcbratney2003}, and where
geostatistics survives it usually enters by assumption: the analyst decides that
residuals should be kriged, selects a variogram family by eye, fixes a lag
binning and an estimator, and adds the correction whether or not it improves
anything. Here we argue that these are not modelling truths to be asserted but
hyperparameters to be selected, and that they should be chosen out of fold
against a null in which no correction is applied, exactly as the number of trees
or the depth of a forest would be. We are candid about what this does and does
not deliver. Regression kriging \citep{hengl2018} and the super learner
\citep{wolpert1992,vanderlaan2007} are established; what is new is that the
geostatistical decisions are made on evidence. And, as our experiment shows,
geostatistics ends up supplying the hypothesis and feature space rather than the
final estimator, since the corrector the data prefers is usually a residual
forest whose features are variogram-scaled basis functions.

The second question is uncertainty. Prediction intervals in soil and ecological
mapping are routinely obtained from kriging variances or quantile heuristics,
and their coverage is almost never checked against the sampling design. Here we
argue that coverage is a population parameter of exactly the same kind as RMSE,
so the design-based apparatus of \citet{wadoux2021} applies to it unchanged and
without further assumptions. This is not a technical refinement. It exposes a
failure that model-based checking cannot see, because model-based checking
relies on the very assumption that clustered sampling breaks.

\section{Coverage is a population parameter, exactly as accuracy is}
\label{sec:ht}

Let the domain be a finite population of $N$ cells, $U=\{1,\dots,N\}$, with true
values $z(s_i)$, and let the map be a function $\hat z: U \rightarrow
\mathbb{R}$. Once produced, the map is fixed and so is the true field. The only
randomness admitted is that of the validation sampling design, which is what
makes the assessment free of assumptions about autocorrelation or stationarity
\citep{degruijter2006,brus2011}.

Let $g: U \rightarrow \mathbb{R}$ be any cell-level function with population
mean $\bar G = N^{-1}\sum_{i \in U} g(i)$, and draw a probability sample
$S \subset U$ of size $n$ with inclusion probabilities $\pi_i$. The
Horvitz-Thompson estimator \citep{horvitz1952},
\[
\widehat{\bar G} = \frac{1}{N}\sum_{i \in S} \frac{g(i)}{\pi_i},
\eqnum{1}
\]
is design-unbiased, and the proof is one line: with $I_i = \mathbf{1}\{i \in S\}$
we have $\mathbb{E}[I_i] = \pi_i$ over repeated sampling, so
\[
\mathbb{E}\big[\widehat{\bar G}\big]
= \frac{1}{N}\sum_{i \in U} \frac{\mathbb{E}[I_i]}{\pi_i}\,g(i)
= \frac{1}{N}\sum_{i \in U} g(i) = \bar G .
\eqnum{2}
\]
Under simple random sampling without replacement $\pi_i = n/N$ and Eq.~(1)
collapses to the sample mean.

The point of writing it this generally is that everything else follows from
choosing $g$. Taking $g(i) = (\hat z(s_i)-z(s_i))^2$ gives the population mean
squared error and its design-unbiased estimator,
\[
\widehat{\mathrm{MSE}} = \frac{1}{n}\sum_{i \in S}\big(\hat z(s_i)-z(s_i)\big)^2 ,
\eqnum{3}
\]
which is the estimator \citet{wadoux2021} advocate. We report
$\widehat{\mathrm{RMSE}} = \sqrt{\widehat{\mathrm{MSE}}}$ and note, since it is
usually left unsaid, that the square root is concave, so Jensen's inequality
gives $\mathbb{E}[\widehat{\mathrm{RMSE}}] \le \mathrm{RMSE}$: the root mean
squared error is not exactly design-unbiased even though the mean squared error
is. The bias is $O(n^{-1})$ and negligible here, and it is a property of the
square root rather than of the design. The mean error, obtained with
$g(i) = \hat z(s_i)-z(s_i)$, is exactly unbiased. We also report the model
efficiency coefficient \citep{nash1970}.

Now take $[\hat \ell(s), \hat u(s)]$ to be a nominal $1-\alpha$ prediction
interval and set $g(i) = \mathbf{1}\{\hat \ell(s_i) \le z(s_i) \le \hat u(s_i)\}$.
Its population mean is the prediction-interval coverage probability, and Eq.~(1)
supplies the estimator at once,
\[
\widehat{\mathrm{PICP}} = \frac{1}{n}\sum_{i \in S}
\mathbf{1}\{\hat \ell(s_i) \le z(s_i) \le \hat u(s_i)\}.
\eqnum{4}
\]
Because $g$ is an indicator, $\widehat{\mathrm{PICP}}$ is a sample proportion:
exactly unbiased, with design variance $\mathrm{PICP}(1-\mathrm{PICP})/n$ up to
the finite-population correction, and elementary confidence statements. Among
valid intervals the better one is narrower, measured by the mean width
$\mathrm{MPIW}$, but width without validity is meaningless and the two must be
read together.

We stress how little this required. No new theory was needed, only the
observation that coverage is a population mean. The consequence is nonetheless
substantial: coverage estimated this way is model-free and can therefore convict
an interval method without sharing any of its assumptions.

\section{A correction should enter only on out-of-fold evidence}
\label{sec:correction}

Write the base random forest \citep{breiman2001} as $\hat f$ and decompose
$z(s) = \hat f(s) + r(s)$. The residual field is not noise: a forest fitted on a
finite sample cannot represent every relation, and what it missed remains in
$r$. If $r$ is predictable then a second model $\hat g$ carries information the
map is discarding, which is the logic of regression kriging \citep{hengl2018}.

The question is how much of it to add. For $\hat y = \hat f + \lambda \hat g$
the population risk is a convex parabola in $\lambda$,
\[
\mathcal{R}(\lambda)
= \mathbb{E}[r^2] - 2\lambda\,\mathbb{E}[r\hat g] + \lambda^2\,\mathbb{E}[\hat g^2],
\eqnum{5}
\]
minimised at
\[
\lambda^\star = \frac{\mathbb{E}[r\hat g]}{\mathbb{E}[\hat g^2]},
\qquad
\mathcal{R}(\lambda^\star) = \mathbb{E}[r^2]
- \frac{\mathbb{E}[r\hat g]^2}{\mathbb{E}[\hat g^2]} \le \mathcal{R}(0).
\eqnum{6}
\]
Three things follow. At the optimum the correction cannot hurt, and helps in
proportion to how well $\hat g$ correlates with the residual. Adding the raw
prediction, $\lambda = 1$, is optimal only if $\hat g$ is perfectly calibrated;
a noisy $\hat g$ inflates $\mathbb{E}[\hat g^2]$ and pushes $\lambda^\star$
below one, so shrinkage is the solution of the risk minimisation and not a
safety device bolted on afterwards. And solving
$\mathcal{R}(\lambda) > \mathcal{R}(0)$ shows the correction damages the map
whenever
\[
\lambda > 2\lambda^\star .
\eqnum{7}
\]
Over-correction is therefore a real failure mode, which is what makes the next
point necessary rather than merely prudent.

Estimating $\lambda$ on the data used to fit $\hat g$ fails, and fails in a
predictable direction. A forest interpolates its training data to some degree,
so the in-sample residual $y_i - \hat f(x_i)$ is shrunk toward zero and is
correlated with the noise in $y_i$. A model fitted to those residuals partly
relearns that noise and appears to explain more of $r$ than it can out of
sample, so the selected $\lambda$ is biased upward, which is exactly the regime
of Eq.~(7). The remedy is $K$-fold cross-fitting: fit the base model on
$U \setminus F_k$ and predict $F_k$ to obtain $\tilde f_i$, and likewise obtain
$\tilde g_i$ from the training folds' residuals. Neither depends on $y_i$, so
\[
\widehat{\mathcal{R}}(c,\lambda)
= \sqrt{\frac{1}{n}\sum_{i=1}^{n}\big(y_i - \tilde f_i - \lambda\,\tilde g_i^{(c)}\big)^2}
\eqnum{8}
\]
estimates out-of-sample risk rather than in-sample fit, the same device that
underlies double machine learning \citep{chernozhukov2018}.

The correction is a library of candidates: a residual forest, a residual Cubist
\citep{quinlan1992}, a single-hidden-layer network, ordinary kriging of the
residuals, a radial basis or kernel smoother, their ensemble, and linear or
quadratic recalibration. We take $(c^\star,\lambda^\star)$ to minimise Eq.~(8).
The structural detail that matters is that $\lambda = 0$ belongs to the grid and
reproduces the base forest, so
\[
\widehat{\mathcal{R}}(c^\star,\lambda^\star) \le \widehat{\mathcal{R}}(\cdot\,,0)
\eqnum{9}
\]
by construction: the method cannot select something worse than the plain forest
on the out-of-fold criterion, and a correction is deployed only after showing an
improvement on data it did not see. This is why the estimator degrades to a
random forest, rather than to something worse, when there is nothing to recover.

Two qualifications belong in the open. The minimum of $\widehat{\mathcal{R}}$
over many candidates is optimistically biased, the familiar winner's curse, so
it is not an unbiased estimate of deployed risk and we never report it as one.
That is precisely why selection happens inside the calibration sample while
accuracy is computed on an independent probability sample: the winner's curse
contaminates the selection score, not the design-based estimate.

Finally, a correction is credible only where data support it. With $d_1(s)$ the
distance to the nearest calibration point we gate it,
\[
w(s) = \exp\!\big(-d_1(s)/\ell\big),
\qquad
\hat y(s) = \hat f(s) + \lambda^\star\, w(s)\, \hat g^{(c^\star)}(s),
\eqnum{10}
\]
with $\ell$ tied to the fitted variogram range, since beyond that range
residuals are effectively uncorrelated and an interpolated correction carries no
information about them. The multiplier is itself selected out of fold, including
the value that disables the gate.

\section{The variogram is a set of hyperparameters, not a set of conventions}
\label{sec:geohyper}

The quantities above are selected, but the variogram feeding them is normally a
fixed recipe: a family chosen by habit, a cutoff at half the maximum distance,
a fixed lag count, one estimator, and a range consumed as a constant. For an
argument that geostatistical choices should rest on evidence, that is the last
unexamined assumption, and we remove it.

The family is selected rather than assumed. With the empirical variogram
computed by the robust estimator of \citet{cressie1980}, the exponential,
spherical, Gaussian and Mat\'ern models \citep{matheron1963,journel1978} are
each fitted to the same empirical variogram and the lowest weighted residual sum
of squares is retained. The comparison is legitimate in this direction, and only
in this direction, because binning and pair counts are identical across
candidates, so the criteria are computed on the same quantity. Parsimony is
enforced among nested families: Mat\'ern nests the exponential at smoothness
$\kappa = 1/2$ and approaches the Gaussian as $\kappa \rightarrow \infty$, and
the smoothness is itself fitted, so raw sum of squares hands Mat\'ern a win by
one extra degree of freedom even when the truth is exactly exponential, which we
observed. We therefore keep the simplest family within a tolerance of the best.

The fitted range enters as a working scale rather than a constant, with radial
basis bandwidth $b = m_b\,\rho$ and the multiplier selected out of fold; the
conventional value corresponds to $m_b = 1/2$. The lag binning and the choice
between the robust and the classical \citep{matheron1963} estimator cannot be
compared by variogram sum of squares at all, because changing them changes the
empirical variogram and hence the quantity being fitted. They are selected by
out-of-fold predictive merit through Eq.~(8).

Anisotropy is treated the same way. We estimate a major-axis angle $\theta$ and
a ratio $\varrho \in (0,1]$ and map coordinates into the space where Euclidean
distance is anisotropic distance,
\[
T_{\theta,\varrho}(s) =
\begin{pmatrix}
\sin\theta & \cos\theta \\[2pt]
-\cos\theta/\varrho & \sin\theta/\varrho
\end{pmatrix} s ,
\eqnum{11}
\]
so that $\varrho = 1$ is a pure rotation recovering the isotropic method
exactly. The transform is applied to the radial basis, the kernel smoother and
the kriging correction, since kriging under geometric anisotropy is ordinary
kriging in the transformed space \citep{journel1978}.

Estimating $(\theta,\varrho)$ deserves a warning, because the obvious routes
fail and we tried them. Taking the ratio of the smallest to the largest
directional range is biased by order statistics, since the maximum of noisy
estimates overshoots and the minimum undershoots; it reports anisotropy for
perfectly isotropic fields. Scoring geometries by variogram fit in the
transformed space fails differently, because the transform changes the distance
distribution and hence the binning, so the criteria are not comparable. We
therefore score each candidate geometry by out-of-fold prediction error of a
kernel smoother on the same held-out points, with bandwidth set as a fraction of
the median pairwise distance in the transformed space so that the score responds
to shape and not to rescaling. On simulated fields this returns exact isotropy
for isotropic fields and the correct angle otherwise, and its errors are false
negatives. Even then the geometry is not imposed: the anisotropic variant enters
the search as one more candidate and must beat the isotropic one on Eq.~(8).

\section{Conformal intervals, and why clustering breaks them}
\label{sec:intervals}

Cross-fitting has already produced leakage-free residuals for the deployed
configuration, and these are exactly the conformity scores split conformal
prediction requires \citep{papadopoulos2002,vovk2005,lei2018}. With $q$ the
$\lceil (n+1)(1-\alpha)\rceil$-th smallest $|e_i|$, each cell receives
$[\hat y(s)-q,\ \hat y(s)+q]$.

The guarantee rests on exchangeability, not on distributional form. If
$|e_1|,\dots,|e_n|,|e_{\mathrm{new}}|$ are exchangeable then $|e_{\mathrm{new}}|$
is equally likely to occupy any rank among the $n+1$ values, so
\[
\Pr\big(|e_{\mathrm{new}}| \le q\big) \ge 1-\alpha
\eqnum{12}
\]
with no assumption of Gaussianity, homoscedasticity or stationarity
\citep{vovk2005}. When calibration and target cells are drawn by simple random
sampling from the same population, exchangeability holds and Eq.~(12) is a
statement about the population, which is design-based coverage in the sense of
Eq.~(4).

Clustered sampling breaks it, and does so directionally. The map must predict
everywhere, including far from every cluster, and the residuals it incurs there
are systematically larger than any observed within clusters, because those
locations are extrapolations. The calibration scores are concentrated on
well-supported locations, so $q$ is calibrated on residuals that are too small,
the intervals are too narrow, and coverage falls below nominal. We test two
repairs. Covariate-shift weighting replaces the empirical quantile by a weighted
one with density-ratio weights \citep{tibshirani2019}; its limitation is
structural rather than technical, since reweighting can only redistribute mass
over residuals that were actually observed. Spatial-block conformal calibration
instead changes what the scores are, refitting on spatial blocks and scoring each
held-out block \citep{roberts2017}, so that every residual is generated in an
extrapolation regime resembling the one the map faces.

\section{Case study}
\label{sec:case}

We use the public materials of \citet{wadoux2021}: above-ground live woody
biomass \citep{baccini2012} over a 928 by 1642 km Amazonian rectangle at 1 km
resolution, 1{,}310{,}143 cells, with 28 covariates known over the whole grid.
Calibration samples of $n = 500$ are drawn under the same three designs they
used, simple random, systematic and two-stage clustered random, and each map is
evaluated design-based on an independent simple random validation sample, with
30 repetitions per design and a fixed seed. Comparators are the untuned default
forest exactly as in the reference study, a default Cubist, and a transparency
model we label RF-tuned, the selected base forest with no correction applied,
which separates the effect of tuning the forest from the effect of the
correction. Forests use \texttt{ranger} \citep{wright2017} and all variogram
work uses \texttt{gstat} \citep{pebesma2004}; every forest receives an explicit
seed. The estimator is released as the \texttt{GeoVersa} package for R.
Fig.~\ref{fig:schematic} summarises the procedure.

\begin{figure}[!ht]
\centering
\includegraphics[width=0.8\textwidth]{figures/fig1_schematic.png}
\caption{The procedure. Cross-fitting on the calibration sample produces
leakage-free out-of-fold predictions and residual candidates; the profile,
correction, shrinkage and gate are selected by out-of-fold RMSE and accepted
only if they beat the RF-only baseline; the winners are refitted on the full
sample and the map is scored on an independent probability sample.}
\label{fig:schematic}
\end{figure}

\subsection{The correction improves the map, and the tuning does not}

The estimator lowers the root mean squared error under all three designs
(Table~\ref{tab:accuracy}), winning 28, 30 and 29 of 30 paired repetitions with
Wilcoxon $p$ values of $6.5\times10^{-9}$, $9.3\times10^{-10}$ and
$1.9\times10^{-9}$, and design-based bias moves toward zero in every design. The
RF-tuned row is the one to read carefully: tuning the forest buys between 0.02
and 0.18 \Mgha{}, one to fifteen per cent of what the full method achieves, so
the improvement is the residual correction and not a better forest. Computed
over all 1{,}310{,}143 cells, that is over the population itself rather than a
sample of it, the root mean squared error is 31.655 \Mgha{} for the default
forest and 31.079 for the corrected map, with bias falling from $+0.378$ to
$+0.127$ (Fig.~\ref{fig:maps}).

\begin{table}[!ht]
\centering
\caption{Design-based accuracy by calibration design, 30 repetitions each,
$n=500$. ME and RMSE in \Mgha{}.}
\label{tab:accuracy}
\begin{tabular}{llrrr}
\toprule
Design & Model & ME & RMSE & MEC \\
\midrule
\multirow{4}{*}{Simple random}
 & RF (default) & $+1.061$ & 32.493 & 0.876 \\
 & RF-tuned     & $+1.129$ & 32.475 & 0.876 \\
 & Cubist       & $+1.853$ & 33.264 & 0.870 \\
 & \textbf{This method} & $\mathbf{+0.826}$ & \textbf{31.624} & \textbf{0.882} \\
\midrule
\multirow{4}{*}{Systematic}
 & RF (default) & $+0.664$ & 32.280 & 0.877 \\
 & RF-tuned     & $+0.707$ & 32.264 & 0.878 \\
 & Cubist       & $+1.921$ & 33.171 & 0.870 \\
 & \textbf{This method} & $\mathbf{+0.479}$ & \textbf{31.400} & \textbf{0.884} \\
\midrule
\multirow{4}{*}{Clustered random}
 & RF (default) & $-1.356$ & 35.778 & 0.851 \\
 & RF-tuned     & $-1.252$ & 35.598 & 0.853 \\
 & Cubist       & $+1.771$ & 37.155 & 0.837 \\
 & \textbf{This method} & $\mathbf{-0.787}$ & \textbf{34.538} & \textbf{0.860} \\
\bottomrule
\end{tabular}
\end{table}

\begin{figure}[!ht]
\centering
\includegraphics[width=\textwidth]{figures/fig7_maps.png}
\caption{Observed biomass, the two predicted surfaces and their error fields for
one calibration sample of $n=500$. Panels (a) to (c) share one colour scale and
panels (d) and (e) share one symmetric scale, so that every comparison is visual
as well as numerical. Red circles in (a) mark the calibration points.}
\label{fig:maps}
\end{figure}

\subsection{The data displaces the conventional variogram recipe}

Table~\ref{tab:axes} reports how often the out-of-fold search departs from the
conventional recipe. In two designs of three it is displaced in a majority of
repetitions, so the defaults were not innocuous.

\begin{table}[!ht]
\centering
\caption{Repetitions, out of 30 per design, in which the search departs from the
conventional recipe (exponential family, cutoff at half the maximum distance, 15
lags, robust estimator, bandwidth of half the range, isotropy).}
\label{tab:axes}
\begin{tabular}{lrrrr}
\toprule
Design & Binning or estimator & Range scale & Anisotropy & Recipe kept \\
\midrule
Simple random    &  9 & 7 & 4 & 10 \\
Systematic       &  4 & 5 & 3 & 18 \\
Clustered random & 11 & 5 & 0 & 14 \\
\bottomrule
\end{tabular}
\end{table}

Two contrasts within the clustered design deserve attention. The anisotropic
geometry is accepted in none of its 30 repetitions, while the binning and
estimator are displaced in 11, the highest count of the three designs. Both are
what a design-aware procedure should do. A clustered sample concentrates points
in a few neighbourhoods and carries almost no information about direction at the
scale of the domain, so the criterion refuses a geometry it cannot estimate; the
same clustering severely distorts the distribution of pairwise distances, so a
cutoff at half the maximum distance with 15 equal lags describes that sample
poorly and the criterion replaces it. The mechanism is not a switch that turns
geostatistics on. It discriminates between the parameters a design can support
and those it cannot.

Table~\ref{tab:modes} reports which corrector is selected across all 90
repetitions, and it is also the limit of what we claim. Ordinary kriging wins
outright in three repetitions and a purely geostatistical interpolator in four;
in 68 the corrector is a residual forest which, carrying a shrinkage weight, is
formally a single boosting step. A reader is entitled to ask whether this is
geostatistics at all.

\begin{table}[!ht]
\centering
\caption{Selected corrector across 90 repetitions.}
\label{tab:modes}
\begin{tabular}{lr}
\toprule
Selected corrector & Count \\
\midrule
Residual random forest & 68 \\
Ensemble of residual forest, Cubist and kernel smoother & 15 \\
Ordinary kriging of residuals & 3 \\
Cubist delta & 2 \\
Kernel smoother & 1 \\
Linear recalibration & 1 \\
\bottomrule
\end{tabular}
\end{table}

It is, but in the representation rather than in the estimator. The residual
forest is not fitted on coordinates alone: its features include radial basis
functions whose centres and bandwidth come from the fitted variogram, and from
the anisotropic geometry of Eq.~(11) when that is accepted, and its output is
multiplied by the gate of Eq.~(10) whose length scale is the fitted range.
Remove the variogram and neither those features nor that gate exist. What
Table~\ref{tab:modes} shows is that once spatial structure is encoded this way a
forest exploits it better than a kriging predictor does on these data, which is
a statement about two learners and not evidence that the geostatistical
information was unnecessary.

\subsection{The advantage is not specific to the scoring protocol}

We drew 10 further calibration samples under simple random sampling and scored
each under three additional protocols (Table~\ref{tab:estimands}): a Population
estimand, an independent simple random sample of 20{,}000 cells excluding the
calibration cells; random 10-fold cross-validation; and spatial block
cross-validation at 350 km.

Two qualifications bound what this shows, and following \citet{wadoux2021} we
state them plainly. Only the first two rows are map-accuracy estimands; the
cross-validation protocols are not valid estimators of map accuracy and are
reported for model comparison only. They are, moreover, aligned with the
internal selection criteria of the method, since corrections are accepted on
random-fold out-of-fold risk and the gate is chosen by spatial
cross-validation, which makes them the least independent of the four checks
rather than the most adversarial. The Population row carries the evidential
weight.

\begin{table}[!ht]
\centering
\caption{Comparison against the untuned default forest under four estimands,
simple random calibration design. RMSE in \Mgha{}.}
\label{tab:estimands}
\begin{tabular}{lrrrrr}
\toprule
Estimand & Reps & RF & This method & $\Delta$RMSE & Wins \\
\midrule
\multicolumn{6}{l}{\emph{Design-based map accuracy}} \\
Population (SRS 20{,}000) & 10 & 32.18 & 31.51 & $-0.669$ & 10/10 \\
Design-based (SRS 500)    & 30 & 32.49 & 31.62 & $-0.869$ & 28/30 \\
\midrule
\multicolumn{6}{l}{\emph{Model comparison only, not map accuracy}} \\
Random 10-fold CV         & 10 & 32.24 & 31.12 & $-1.125$ & 10/10 \\
Spatial block CV (350 km) & 10 & 34.72 & 33.45 & $-1.275$ & 10/10 \\
\bottomrule
\end{tabular}
\end{table}

\begin{figure}[!ht]
\centering
\includegraphics[width=0.88\textwidth]{figures/fig4_protocols.png}
\caption{Comparison under four validation estimands, simple random calibration
design. Population and Design-based estimate map accuracy; the two
cross-validation estimands are reported for model comparison only and are not
valid map-accuracy estimates. The cross-validation bars train on subsets of the
calibration sample, so their absolute heights are not comparable with the first
two.}
\label{fig:protocols}
\end{figure}

The method wins all 10 draws under each additional protocol. Three cautions. The
one-sided paired Wilcoxon $p$ at 10 repetitions cannot fall below
$2^{-10} \approx 0.00098$, the value attained, so it certifies that all ten
draws favour the method and nothing stronger. These are 10 draws scored three
ways rather than 30 independent repetitions, and results within a draw are
dependent, correlating at $r = 0.59$. The cross-validation rows train on strict
subsets, so their absolute values are not commensurable with the design-based
rows. We do not read the ordering of the four differences as a gradient:
adjacent protocols are not statistically distinguishable ($p = 0.36$), and the
ordering is not one of spatial stress, since random 10-fold cross-validation is
the least spatially demanding of the four. Buffered leave-one-out
cross-validation was not run, requiring of order $10^3$ sequential fits per
repetition; it would be the most adversarial estimand available and is left to
future work.

\subsection{Design-based assessment convicts the intervals}

Table~\ref{tab:coverage} contains the finding we consider most consequential.
Under representative sampling the plain conformal intervals are calibrated,
0.896 and 0.907 against a nominal 0.90, and about 17 per cent narrower than the
quantile regression forest baseline \citep{meinshausen2006}, which over-covers
at 0.937. Under clustered sampling they under-cover at 0.840, and
covariate-shift weighting barely helps at 0.850, for the structural reason given
in Section~\ref{sec:intervals}. The spatial-block calibration restores validity
at 0.946 at the cost of intervals about 40 per cent wider.

\begin{table}[!ht]
\centering
\caption{Design-based interval coverage at a nominal 0.90 level and mean width in
\Mgha{}.}
\label{tab:coverage}
\begin{tabular}{llrr}
\toprule
Design & Interval & PICP & MPIW \\
\midrule
\multirow{2}{*}{Simple random}
 & Conformal & 0.896 & 100 \\
 & Quantile regression forest & 0.937 & 123 \\
\midrule
Systematic & Conformal & 0.907 & 104 \\
\midrule
\multirow{3}{*}{Clustered random}
 & Conformal, plain & 0.840 & 111 \\
 & Covariate-shift weighted & 0.850 & 114 \\
 & \textbf{Spatial-block conformal} & \textbf{0.946} & 145 \\
\bottomrule
\end{tabular}
\end{table}

What matters is not any single number but that the failure is visible at all. A
model-based coverage check would have relied on the same exchangeability that
clustering breaks, and would have reported these intervals as adequate.
Estimating coverage as a population parameter on a probability sample is what
convicts them, and it does so without sharing any assumption with the method
under test.

\begin{figure}[!ht]
\centering
\includegraphics[width=\textwidth]{figures/fig6_uncertainty.png}
\caption{Design-based interval assessment. Left, coverage with the nominal 0.90
level marked; right, mean interval width. The plain conformal calibration is
valid and sharpest under representative sampling but under-covers under
clustering, and only the spatial-block calibration restores coverage, at the
cost of width.}
\label{fig:uncertainty}
\end{figure}

\section{Discussion}

Our first argument is about where decisions are made, not about a new learner.
The ingredients are established and we say so. What the experiment supports is
narrower and, we think, more useful: when geostatistical decisions are moved
into the hyperparameter space and validated out of fold, the map improves
reliably under the design-based yardstick, and it does so without the analyst
judging a variogram by eye or deciding in advance whether kriging is warranted.
The RF-tuned control localises the gain in the correction, Eq.~(9) bounds the
downside by the forest itself, and Table~\ref{tab:axes} shows the freedom is used
and used discriminatingly.

We are equally clear about the boundary, because overstating it would repeat the
error \citet{wadoux2021} criticised. The selected corrector is usually a
residual forest and ordinary kriging wins outright in three of 90 repetitions.
We do not claim that kriging has returned to the centre of digital soil mapping.
We claim that the geostatistical questions, at what scale, in which geometry,
with how much confidence and whether at all, are posed explicitly and answered
by evidence, and that the answers change what the model does even when the final
corrector is a forest.

The second argument is the more consequential. Extending design-based assessment
from error to coverage took three lines in Section~\ref{sec:ht}, but it changes
what can be checked. Intervals that look entirely reasonable can be badly
miscalibrated for the population under clustered sampling, which is the same
regime in which \citet{wadoux2021} found point-accuracy estimators to fail. The
recommendation follows directly: assess coverage design-based, and use a
spatial-block calibration when clustering is suspected, accepting wider but
valid intervals.

Several limitations qualify all of this. The point gain is incremental, of order
0.9 to 1.2 \Mgha{}. Opening the geostatistical hyperparameters leaves accuracy
essentially unchanged relative to the fixed recipe, so what is demonstrated is
that the choices are made on evidence and are shown to matter to the fitted
model, not that they lower error on these data. Anisotropy is accepted in about
one repetition in five under representative sampling and never under clustering,
and the model is geometric and global, one angle and ratio for the whole domain.
These results come from a single, large and widely used dataset at one sample
size.

\section{Conclusion}

Geostatistics need not be assumed into a map nor exiled from it. Its parameters,
the variogram family, the range, the binning, the estimator and the anisotropy,
can be treated as hyperparameters and asked to prove themselves out of fold, and
when they are, the conventional recipe is displaced in most repetitions and the
procedure declines the parameters the sampling design cannot support. On these
data the corrector the model finally prefers is a forest, but it is a forest
that sees the landscape through a variogram.

Coverage, meanwhile, is a population parameter of exactly the same kind as RMSE,
and inherits the same design-unbiased estimator. This is worth saying because it
is easy to miss and consequential when missed: under clustered sampling
conformal intervals under-cover, no model-based check can reveal it, and a
spatial-block calibration repairs it. Map uncertainty, like map accuracy, must
be assessed design-based.

\section*{Declaration of generative AI in the writing process}

During the preparation of this work the author used Grammarly for spelling and
grammar checking, and an AI assistant (Claude, Anthropic) to draft manuscript
text, implement analysis code and produce figures, under the author's direction.
All study design, methodological decisions, interpretation of results and
verification of the analysis are the author's own. After using these tools the
author reviewed and edited the content as needed and takes full responsibility
for the content of the publication.

\section*{CRediT authorship contribution statement}

\textbf{Hugo Rodrigues}: Conceptualization, Methodology, Software, Formal
analysis, Investigation, Data curation, Writing, Visualization.

\section*{Declaration of competing interest}

The author declares no known competing financial interests or personal
relationships that could have appeared to influence the work reported in this
paper.

\section*{Data availability}

The validation data are the public materials of \citet{wadoux2021}, available at
github.com/AlexandreWadoux/SpatialValidation, and derive from
\citet{baccini2012}. The estimator is released as the \texttt{GeoVersa} package
for R on CRAN, and the analysis code is available in the project repository.

\bibliographystyle{elsarticle-harv}
\bibliography{references}

\end{document}
